\section{Experimental design}\label{sec:setup}

\paragraph{Research questions}
RQ1: Which explainers localize the attacked loop at the alarm, and does control-loop structure matter beyond grouping per se? RQ2: How does accuracy evolve after the alarm, and does \TEA{} help? RQ3: Which explainers deliver correct localizations before a deadline once detection delay and latency are accounted for? RQ4: How robust are they to sensor noise and a decoy manipulation? RQ5: How much does a random window split inflate detection results? Hypotheses H1--H6 about \LX{} and \TEA{} and the decision rule (two-sided tests, Holm-corrected $p<0.05$ with the predicted sign) were fixed in the project's design document before the main runs (Supplementary Section~S1). The comparison with PI-CAX was added afterwards and is reported as a post hoc comparison; PI-CAX's design and constants were frozen before it was first run on either benchmark.

\subsection{Benchmarks}
\paragraph{HAI 21.03} HAI 21.03 \citep{shin2021hai} is a hardware-in-the-loop testbed coupling a boiler (P1), a turbine (P2) and a water-treatment process (P3) with a simulated power plant (P4). It records 79 tags at 1\,Hz: three normal training files (9.5 days) and five test files with 50 attacks (25 single, 25 combined). Its documentation lists the target controller(s) and manipulated points of every attack.
\begin{itemize}
\item \emph{Ground truth.} We transcribed this table and matched every entry to the label runs in the data (all start times agree within 2\,min); three point names were mapped to their 21.03 columns. An independent transcription of the same table, produced separately within the project, agrees on the target loops and points of \gtcrossAgree{} of \gtcrossN{} attacks. The one disagreement (A401) is a row continuing across a page break; we keep its second scenario, as for every A4xx attack, and excluding A401 changes no reported HR@1 by more than \aFourZeroOneMaxDiff.
\item \emph{Loop map.} Tags were assigned to loop groups from the tag descriptions before any test data was examined (Supplementary Section~S3). Twenty-two tags are constant in normal training and removed, leaving $d=\HAId$ tags in $L=\HAIL$ groups.
\item \emph{Hidden setpoints.} Manipulated setpoints are largely invisible in the recorded data: of \spN{} (attack, setpoint-tag) targets, \spConstant{} are exactly constant during the attack and only \spVisible{} vary beyond their range in the preceding 10\,min. Tag-level localization of those points is impossible from the data; loop-level localization is not.
\end{itemize}

\paragraph{PCL-Sim} The simulator models a supply tank, a transfer valve, a heated tank, a centrifugal pump and a discharge valve under five PI loops (tank-1 level, transfer flow, tank-2 level via pump speed, tank-2 temperature, discharge pressure), with actuator dynamics, hydraulic and thermal coupling, sensor noise, random-walk disturbances and an operator schedule that changes setpoints every 15--45\,min. It exposes $d=22$ tags in $L=6$ groups (five loops and a disturbance group). Seven attack primitives mirror HAI (SP, CV offset, PV bias, PV freeze, CV+PV-freeze, SP+PV-freeze, short CV pulses), applied to each loop at two magnitudes and four test seeds, for \SimNattacks{} one-hour episodes. Controllers and historian see the manipulated values while the true state is integrated separately; the hazard deadline $D^{\mathrm{haz}}$ is the time to the first violation of a safety limit (overflow or dry-run, temperature $>85\,^\circ$C, pressure $>8.5$\,bar). \SimNhaz{} attacks are hazardous (median $D^{\mathrm{haz}}=\SimHazMedian$\,s, range \SimHazMin--\SimHazMax\,s). Twelve 6-hour normal runs with disjoint seeds are used for training and validation.

\subsection{Detectors and leakage control}
Four forecaster families are trained on normal data only, with three seeds each and window $w=32$: a linear autoregressive model (closed-form ridge), an MLP, a dilated 1D-CNN \citep[cf.][]{kravchik2018cnn} and a GRU. Normal files are split chronologically (80\,\% train, 20\,\% validation); standardisation comes from training data, and $\sigma_j$ and $\theta$ ($q=0.999$, $\alpha=0.1$) from validation data only. No test labels, files or seeds are used for any fitting or hyperparameter choice.

\subsection{Explainers}
All explainers attribute the same score $e(Z_t)$; tag-level attributions are summed per loop.
\begin{itemize}
\item \emph{Tag-level baselines:} raw squared residuals (RE; the ``MSE'' ranking of \citealp{fung2024attributions}), saliency (SM; \citealp{simonyan2014saliency}), integrated gradients from the reference (IG, 32 steps; \citealp{sundararajan2017ig}), KernelSHAP \citep{lundberg2017shap} and LIME \citep{ribeiro2016lime} with $2^L$ (``eq'') or 1024 samples, and the ensemble of \citet{fung2024attributions} (RE + SM + LIME in place of LEMNA, $\beta=2.5$ for actuators).
\item \emph{Loop-structured:} \LX, \LX-bg (\LX{} with an interventional background of 8 training windows) and PI-CAX (\cref{sec:picax}).
\item \emph{Structure controls:} exact Shapley over $L$ correlation clusters (CLX) or a random partition with the loop group sizes (RPX).
\item \emph{ShaTS-sensor} \citep{franco2026shats}, with $m=20\,|G|$ coalitions per group; its cost ($\approx2\cdot20\,d^2$ evaluations) restricts it to two detectors at the alarm.
\end{itemize}
Every method is run single-shot and with \TEA.

\subsection{Protocol, metrics and statistics}
Following \citet{fung2024attributions}, explanations are computed at a best-guess timing (oracle window starting at onset) and at practical timing: from the first alarm every $\Delta=5$\,s for 180\,s.
\begin{itemize}
\item \emph{Detection:} point-wise precision, recall and F1 without point adjustment \citep{kim2022rigorous}, event recall (alarm within the attack or 60\,s after it), median delay and false alarms per hour.
\item \emph{Localization:} HR@$k$ (a target loop among the top $k$) and normalised rank, against a uniformly random ranking as chance reference.
\item \emph{Faithfulness:} deletion AUC, restoring loops in ranked order to the training mean (lower is better).
\item \emph{Stability:} top-$k$ agreement between consecutive explanations, and Jaccard agreement under Gaussian input noise ($\sigma=0.1$).
\item \emph{Cost and timeliness:} median wall-clock latency on one CPU thread, detector evaluations per explanation, DCL@$k(D)$ and the in-time rate with $t_{\mathrm{act}}=30$\,s.
\end{itemize}
The unit of analysis is the attack: seeds are averaged per (attack, detector), and these pairs are the paired units. We report means with 95\,\% bootstrap confidence intervals and test with a Friedman omnibus and two-sided Wilcoxon signed-rank tests against a reference configuration, Holm-corrected \citep{demsar2006statistical,holm1979simple}, with matched-pairs rank-biserial effect sizes $r_{\mathrm{rb}}$. Experiments ran with Python~3.13 and PyTorch~2.15 on an \CPUmodel{}; latencies use one thread (\texttt{torch.set\_num\_threads(1)}).
